\chapter{The Visco-Elastic Derivation of Cosmic Redshift}
\label{app:cosmic_redshift}
\addtocontents{toc}{\protect\setcounter{tocdepth}{0}}

\textbf{Abstract:} Standard cosmology derives redshift ($z$) from the kinematic stretching of the metric. Tension Field Theory (TFT) derives it strictly from the visco-elastic drag of a wave propagating through a dissipative medium. This appendix derives the raw kinematic redshift coordinate distance.

\section{The Energy Decay Function}
As a photon propagates through the Spacetime Gravitational Floor (SGF), it experiences a constant fractional energy loss per unit distance due to Lattice Hysteresis (detailed in Section 3.3). This dissipated transverse energy is permanently deposited into the lattice as longitudinal phonon heat, sustaining the 2.725 K CMB. This defines the macroscopic Lattice Damping Rate ($\Gamma_D$). The differential equation for this energy loss is:
\begin{equation}
    \frac{dE}{dx} = -\frac{\Gamma_D}{c_E} E
\end{equation}

Integrating this over a propagation distance ($x$) yields a compounding exponential decay function:
\begin{equation}
    E_{obs} = E_{emit} e^{-\frac{\Gamma_D x}{c_E}}
\end{equation}

\section{Deriving the Redshift Factor ($1+z$)}
Because the energy of a photon is inversely proportional to its wavelength ($E \propto 1/\lambda$), we can substitute wavelength into the decay function:
\begin{equation}
    \frac{1}{\lambda_{obs}} = \frac{1}{\lambda_{emit}} e^{-\frac{\Gamma_D x}{c_E}}
\end{equation}
\begin{equation}
    \lambda_{obs} = \lambda_{emit} e^{\frac{\Gamma_D x}{c_E}}
\end{equation}

The standard astronomical definition of redshift is $1+z = \lambda_{obs}/\lambda_{emit}$. Substituting our derived wavelength:
\begin{equation}
    1+z = e^{\frac{\Gamma_D x}{c_E}}
\end{equation}

\section{Resolving the Hubble Law}
For relatively short, local cosmic distances, the Taylor series expansion of the exponential function ($e^y \approx 1+y$) perfectly reduces to a linear approximation:
\begin{equation}
    1+z \approx 1 + \frac{\Gamma_D x}{c_E}
\end{equation}
\begin{equation}
    z \approx \frac{\Gamma_D}{c_E} x
\end{equation}

\begin{itemize}
    \item This matches the exact linear format of Hubble's Law ($v = H_0 d$).
    \item By setting $\Gamma_D \equiv H_0$, TFT proves that the Hubble constant is not a velocity of metric expansion, but the literal hydrodynamic drag coefficient of the vacuum.
\end{itemize}

Furthermore, at extreme distances (such as Type Ia Supernovae), the higher-order terms of the exponential expansion can no longer be ignored. Substituting $t = x/c_E$, the second-order expansion yields:
\begin{equation}
    z \approx \Gamma_D t + \frac{(\Gamma_D t)^2}{2}
\end{equation}

This mathematically proves that the inherent exponential nature of the decay natively produces a non-linear, upward-curving redshift. The $+\frac{(\Gamma_D t)^2}{2}$ term exactly mimics the ``dimming'' that standard cosmology incorrectly attributes to an accelerating metric (``Dark Energy'').